% --- Archon physics formalization source begin ---
% archon:physics
% archon:covers PhyXMiniProblems/problem_phyx_mini_0363.lean
% archon:source-report reports/phyx_mini/problem_phyx_mini_0363.source.json

\chapter{Physics problem phyx\_mini\_0363}
\label{ch:PhyXMiniProblems_problem_phyx_mini_0363}

\paragraph{Problem source.}
\#\# Physical scenario

0.020 mol of gas undergoes the process shown in the figure.

\#\# Question

What is the final temperature?

\#\# Answer choices

- A: A: 458\textasciicircum{}\textbackslash{}circ C

- B: B: 756\textasciicircum{}\textbackslash{}circ C

- C: C: 356\textasciicircum{}\textbackslash{}circ C

- D: D: 584\textasciicircum{}\textbackslash{}circ C

\#\# Figure caption (auxiliary; use the image as primary evidence)

The image is a graph with the x-axis labeled as \textbackslash{}( V \textbackslash{}, (\textbackslash{}text\{cm\}\textasciicircum{}3) \textbackslash{}) and the y-axis labeled as \textbackslash{}( p \textbackslash{}, (\textbackslash{}text\{atm\}) \textbackslash{}). The graph illustrates a downward curving line representing a process between two points, labeled 1 and 2.

- Point 1 is at a higher pressure and lower volume.
- Point 2 is at a lower pressure and higher volume.
- An arrow indicates the direction from point 1 to point 2, suggesting a decrease in pressure and an increase in volume.
- The x-axis is marked at 0 and \textbackslash{}( V\_2 \textbackslash{}).
- The y-axis ranges from 0 to 3.

The curve illustrates an inverse relationship between the pressure and volume, likely representing an isothermal or adiabatic process in thermodynamics.

\#\# Dataset metadata

Ideal Gases and Kinetic Theory; Physical Model Grounding Reasoning, Multi-Formula Reasoning

\paragraph{Recorded answer/context.}
A: A: 458\textasciicircum{}\textbackslash{}circ C

\paragraph{Figure/image path.}
/root/proposal\_for\_physic/science-mango/phyx\_mini\_run/phyx\_data/test\_image/363.png

\paragraph{Formalization target.}
create a compiling Lean file with sorry bodies at `PhyXMiniProblems/problem\_phyx\_mini\_0363.lean`. The Lean declarations must preserve the physical quantities, dimensions or dimensional roles, figure labels, governing-law hypotheses, and final relation expressed by this problem.
Use Mathlib/Physlib names found through LeanExplore where available. If a physics API is missing, introduce faithful local abstractions rather than scalar placeholder aliases.

\begin{theorem}[Physics formalization target]
\label{thm:physics:phyx_mini_0363:target}
\lean{PhyXMiniProblems.ProblemPhyXMini0363.finalTemperatureRoundsTo458C}
\uses{def:physics:phyx-mini-0363:phyxminiproblems-problemphyxmini0363-pressureinatmospheres, def:physics:phyx-mini-0363:phyxminiproblems-problemphyxmini0363-volumeincubiccentimeters, def:physics:phyx-mini-0363:phyxminiproblems-problemphyxmini0363-temperatureinkelvins, def:physics:phyx-mini-0363:phyxminiproblems-problemphyxmini0363-directedgasprocess, def:physics:phyx-mini-0363:phyxminiproblems-problemphyxmini0363-haspositivereadouts, def:physics:phyx-mini-0363:phyxminiproblems-problemphyxmini0363-satisfiesidealgaslaw, def:physics:phyx-mini-0363:phyxminiproblems-problemphyxmini0363-followsinversepvcurve, def:physics:phyx-mini-0363:phyxminiproblems-problemphyxmini0363-roundstonearestdegreecelsius}
For the directed inverse-\(pV\) process with \(n=0.020\) mole,
\(R=82.057\,\mathrm{atm\,cm^3/(mol\,K)}\), and first-state readouts
\(p_1=3\,\mathrm{atm}\), \(V_1=400\,\mathrm{cm^3}\), the final absolute
temperature rounds to \(458^\circ\mathrm C\).
\end{theorem}
\begin{proof}
The inverse-curve hypothesis makes \(p_1V_1=p_2V_2\).  Applying the ideal-gas
law at both states with the same positive \(nR\) therefore gives
\(T_2=T_1\).  At state one,
\[
 T_1=\frac{3\cdot400}{(20/1000)(82057/1000)}
     =\frac{60000000}{82057}\ {\rm K}.
\]
Subtract \(273.15\) to obtain the Celsius readout.  Direct rational
comparison shows that its distance from \(458\) is strictly below one half,
which is precisely the declared nearest-degree predicate.
\end{proof}

% --- Archon declaration topology begin ---
\section{Declaration topology}
\begin{definition}[Volume Quantity]
  \label{def:physics:phyx-mini-0363:phyxminiproblems-problemphyxmini0363-volumequantity}
  \lean{PhyXMiniProblems.ProblemPhyXMini0363.VolumeQuantity}
  A physical volume, tagged with dimension length cubed and a nonnegative value
\end{definition}
\begin{proof}
  This is the declared model definition of Volume Quantity.
\end{proof}
\begin{definition}[Pressure Quantity]
  \label{def:physics:phyx-mini-0363:phyxminiproblems-problemphyxmini0363-pressurequantity}
  \lean{PhyXMiniProblems.ProblemPhyXMini0363.PressureQuantity}
  Physical pressure, using Physlib's dimensionful pressure type
\end{definition}
\begin{proof}
  This is the declared model definition of Pressure Quantity.
\end{proof}
\begin{definition}[Temperature Quantity]
  \label{def:physics:phyx-mini-0363:phyxminiproblems-problemphyxmini0363-temperaturequantity}
  \lean{PhyXMiniProblems.ProblemPhyXMini0363.TemperatureQuantity}
  Absolute thermodynamic temperature
\end{definition}
\begin{proof}
  This is the declared model definition of Temperature Quantity.
\end{proof}
\begin{definition}[pressure In Pascals]
  \label{def:physics:phyx-mini-0363:phyxminiproblems-problemphyxmini0363-pressureinpascals}
  \lean{PhyXMiniProblems.ProblemPhyXMini0363.pressureInPascals}
  \uses{def:physics:phyx-mini-0363:phyxminiproblems-problemphyxmini0363-pressurequantity}
  Coherent-SI pressure readout in pascals
\end{definition}
\begin{proof}
  This is the declared model definition of pressure In Pascals.
\end{proof}
\begin{definition}[pressure In Atmospheres]
  \label{def:physics:phyx-mini-0363:phyxminiproblems-problemphyxmini0363-pressureinatmospheres}
  \lean{PhyXMiniProblems.ProblemPhyXMini0363.pressureInAtmospheres}
  \uses{def:physics:phyx-mini-0363:phyxminiproblems-problemphyxmini0363-pressurequantity, def:physics:phyx-mini-0363:phyxminiproblems-problemphyxmini0363-pressureinpascals}
  Pressure readout in standard atmospheres
\end{definition}
\begin{proof}
  This is the declared model definition of pressure In Atmospheres.
\end{proof}
\begin{definition}[volume In Cubic Centimeters]
  \label{def:physics:phyx-mini-0363:phyxminiproblems-problemphyxmini0363-volumeincubiccentimeters}
  \lean{PhyXMiniProblems.ProblemPhyXMini0363.volumeInCubicCentimeters}
  \uses{def:physics:phyx-mini-0363:phyxminiproblems-problemphyxmini0363-volumequantity}
  Volume readout in cubic centimetres; 1 m³ = 10⁶ cm³
\end{definition}
\begin{proof}
  This is the declared model definition of volume In Cubic Centimeters.
\end{proof}
\begin{definition}[temperature In Kelvins]
  \label{def:physics:phyx-mini-0363:phyxminiproblems-problemphyxmini0363-temperatureinkelvins}
  \lean{PhyXMiniProblems.ProblemPhyXMini0363.temperatureInKelvins}
  \uses{def:physics:phyx-mini-0363:phyxminiproblems-problemphyxmini0363-temperaturequantity}
  Convert the stored absolute-temperature scale to a kelvin readout
\end{definition}
\begin{proof}
  This is the declared model definition of temperature In Kelvins.
\end{proof}
\begin{definition}[Gas State]
  \label{def:physics:phyx-mini-0363:phyxminiproblems-problemphyxmini0363-gasstate}
  \lean{PhyXMiniProblems.ProblemPhyXMini0363.GasState}
  \uses{def:physics:phyx-mini-0363:phyxminiproblems-problemphyxmini0363-volumequantity, def:physics:phyx-mini-0363:phyxminiproblems-problemphyxmini0363-pressurequantity, def:physics:phyx-mini-0363:phyxminiproblems-problemphyxmini0363-temperaturequantity}
  A thermodynamic state with physical pressure, volume, and absolute temperature
\end{definition}
\begin{proof}
  This is the declared model definition of Gas State.
\end{proof}
\begin{definition}[Directed Gas Process]
  \label{def:physics:phyx-mini-0363:phyxminiproblems-problemphyxmini0363-directedgasprocess}
  \lean{PhyXMiniProblems.ProblemPhyXMini0363.DirectedGasProcess}
  \uses{def:physics:phyx-mini-0363:phyxminiproblems-problemphyxmini0363-gasstate}
  The directed process arrow in the figure, from point 1 to point 2
\end{definition}
\begin{proof}
  This is the declared model definition of Directed Gas Process.
\end{proof}
\begin{definition}[Has Positive Readouts]
  \label{def:physics:phyx-mini-0363:phyxminiproblems-problemphyxmini0363-haspositivereadouts}
  \lean{PhyXMiniProblems.ProblemPhyXMini0363.HasPositiveReadouts}
  \uses{def:physics:phyx-mini-0363:phyxminiproblems-problemphyxmini0363-pressureinatmospheres, def:physics:phyx-mini-0363:phyxminiproblems-problemphyxmini0363-volumeincubiccentimeters, def:physics:phyx-mini-0363:phyxminiproblems-problemphyxmini0363-temperatureinkelvins, def:physics:phyx-mini-0363:phyxminiproblems-problemphyxmini0363-gasstate}
  Positivity of the pressure, volume, and kelvin readouts of a state
\end{definition}
\begin{proof}
  This is the declared model definition of Has Positive Readouts.
\end{proof}
\begin{definition}[Satisfies Ideal Gas Law]
  \label{def:physics:phyx-mini-0363:phyxminiproblems-problemphyxmini0363-satisfiesidealgaslaw}
  \lean{PhyXMiniProblems.ProblemPhyXMini0363.SatisfiesIdealGasLaw}
  \uses{def:physics:phyx-mini-0363:phyxminiproblems-problemphyxmini0363-pressureinatmospheres, def:physics:phyx-mini-0363:phyxminiproblems-problemphyxmini0363-volumeincubiccentimeters, def:physics:phyx-mini-0363:phyxminiproblems-problemphyxmini0363-temperatureinkelvins, def:physics:phyx-mini-0363:phyxminiproblems-problemphyxmini0363-gasstate}
  The ideal-gas law written in the chart units atm, cm³, mol, and K
\end{definition}
\begin{proof}
  This is the declared model definition of Satisfies Ideal Gas Law.
\end{proof}
\begin{definition}[Follows Inverse PVcurve]
  \label{def:physics:phyx-mini-0363:phyxminiproblems-problemphyxmini0363-followsinversepvcurve}
  \lean{PhyXMiniProblems.ProblemPhyXMini0363.FollowsInversePVcurve}
  \uses{def:physics:phyx-mini-0363:phyxminiproblems-problemphyxmini0363-pressureinatmospheres, def:physics:phyx-mini-0363:phyxminiproblems-problemphyxmini0363-volumeincubiccentimeters, def:physics:phyx-mini-0363:phyxminiproblems-problemphyxmini0363-directedgasprocess}
  Model the pictured inverse p-V curve by a constant pressure-volume product between its labeled endpoints
\end{definition}
\begin{proof}
  This is the declared model definition of Follows Inverse PVcurve.
\end{proof}
\begin{definition}[kelvin To Celsius]
  \label{def:physics:phyx-mini-0363:phyxminiproblems-problemphyxmini0363-kelvintocelsius}
  \lean{PhyXMiniProblems.ProblemPhyXMini0363.kelvinToCelsius}
  Convert an absolute-temperature readout in kelvin to degrees Celsius
\end{definition}
\begin{proof}
  This is the declared model definition of kelvin To Celsius.
\end{proof}
\begin{definition}[Rounds To Nearest Degree Celsius]
  \label{def:physics:phyx-mini-0363:phyxminiproblems-problemphyxmini0363-roundstonearestdegreecelsius}
  \lean{PhyXMiniProblems.ProblemPhyXMini0363.RoundsToNearestDegreeCelsius}
  \uses{def:physics:phyx-mini-0363:phyxminiproblems-problemphyxmini0363-kelvintocelsius}
  A Celsius value rounds to the displayed whole-degree answer
\end{definition}
\begin{proof}
  This is the declared model definition of Rounds To Nearest Degree Celsius.
\end{proof}
% --- Archon declaration topology end ---
% --- Archon physics formalization source end ---
